% 第四章：理论、算法与实现映射

\section{算法设计与实现映射} \label{sec:theory}

\subsection{系统模型}

系统有 $m$ 个同构核 $\Lambda=\{\lambda_1,\dots,\lambda_m\}$，与 $\eta$ 级包容式缓存 $L=\{L_1,\dots,L_\eta\}$：$L_1$ 私有，$L_2$ 由相邻 $\kappa$ 个核共享，$L_3$ 全芯片共享。任务集 $\Gamma=\{\tau_1,\dots,\tau_n\}$，$\tau_i=\{G_i,T_i,P_i,w_i\}$，$w_i=\sum_{u\in V_i}C_u$。作业 $u_{i,j,v}$ 是节点 $u_{i,j}$ 的第 $v$ 次实例。

记核 $\lambda_k$ 上已排定工作的完成时刻为 $\ell_k$（负载），作业的基准执行时间为 $\hat{C}$，预测执行时间为 $\hat{E}(\cdot)$，放置决策要解的问题是
\begin{equation}
\min_{\alpha:\;u\mapsto\lambda} \max_k \ell_k,
\end{equation}
它是带缓存依赖的 $P\,\|\,C_{\max}$ 变体：加工时间不是常数，而取决于同一作业此前在该核上的执行历史。该问题没有已知的多项式最优解，工程上只能用启发式。

\subsection{重用距离、CRP 与两条规则}

第 $x$ 级缓存上的重用距离用时间近似（\autoref{eq:reuse}），CRP 把重用距离映射为千分比 $\rho\in[0,1000]$：
\begin{equation}\label{eq:reuse}
r(u_j,L_x)=g\big(t(u_{j,v})-t(u_{j,v-1})\big),\qquad
\hat{E}(u_j,\lambda_k)=\Big(1-\frac{\rho\big(r(u_j,L_x)\big)}{1000}\Big)\hat{C}.
\end{equation}

两条分配规则：

\textbf{MSF}：在候选核集合上取预测加速比最大者，
$\lambda^\ast=\arg\max_{\lambda_k\in\Lambda^\ast}\big(\hat{C}-\hat{E}(u_j,\lambda_k)\big)$。

\textbf{LCIF}：加速比并列时，取对其他作业缓存收益破坏最小者（\autoref{eq:lcif}），
\begin{equation}\label{eq:lcif}
\mathrm{Imp}(u_j,\lambda_k)=\sum_{u_x\in H(\lambda_k)}\Big(\hat{E}(u_x;\lambda_k,H)-\hat{E}(u_x;\lambda_k,H')\Big).
\end{equation}

\autoref{alg:ajlr} 给出在线流程，其中 $\Lambda^\ast$ 是候选核集合，$Q_{sched}$ 是本调度点待分派的作业。

\begin{algorithm}[H]
\caption{Yat-CASched 在线分配}
\label{alg:ajlr}
\KwIn{就绪作业 $Q_{ready}$，候选核 $\Lambda^\ast$，历史表 $H$，CRP}
$Q_{sched}\leftarrow \mathrm{sort}(Q_{ready})[0:|\Lambda^\ast|]$\;
\For{$u_j\in Q_{sched}$}{
  \For{$\lambda_k\in\Lambda^\ast$}{
    $S(u_j,\lambda_k)\leftarrow \hat{C}_j-\hat{E}(u_j,\lambda_k;H,\mathrm{CRP})$\;
  }
}
\While{$Q_{sched}\neq\emptyset$}{
  $(u_j,\Lambda^{\neg})\leftarrow\arg\max\{S(u_j,\lambda_k)\}$\;
  \eIf{$|\Lambda^{\neg}|=1$}{$\lambda_k\leftarrow\Lambda^{\neg}[0]$\;}{
    $\lambda_k\leftarrow\arg\min_{\lambda\in\Lambda^{\neg}}\mathrm{Imp}(u_j,\lambda)$\;}
  $H(\lambda_k)\leftarrow H(\lambda_k)\cup\{u_j\}$，$Q_{sched}\leftarrow Q_{sched}\setminus\{u_j\}$\;
}
\end{algorithm}

单次决策的代价是 $O(|\Lambda^\ast|\cdot \bar h)$，$\bar h$ 是历史表的平均扫描长度；一轮批量分配至多 $O(m^2\bar h)$。历史表长度必须有硬上限，否则"在线"会退化为"离线"。

\subsection{设计到实现的映射}

概念与内核实现的对应关系见\autoref{tab:mapping}。"实现里怎么算"这一列不是伪代码，而是
\texttt{yat\_casched.c} 里的函数名——论文与代码之间靠这张表对齐，而不是靠叙述。

\begin{table}[H]
\centering
\caption{概念 $\to$ 内核实现}
\label{tab:mapping}
\small
\begin{tabular}{@{}p{2.6cm}p{4.6cm}p{6.4cm}@{}}
\toprule
概念 & 内核实现 & 关键取舍 \\
\midrule
作业与身份 & \texttt{struct yat\_job}，\texttt{job\_id} 与 \texttt{task\_struct} 同寿 & 作业身份必须跨唤醒保持，否则重用距离无从累计 \\
重用距离 & \texttt{yat\_calculate\_recency()}：按核累计执行区间 & 用时间近似访问序列，省掉逐访问记账 \\
CRP & \texttt{yat\_get\_crp\_ratio()}：千分比定点 & 定点避免浮点，且便于在 tracepoint 里原样打印 \\
MSF & \texttt{yat\_find\_best\_core\_allowed()}：benefit 最大 & 只在"可运行且有历史"的核上算，避免全表扫描 \\
破平 & 同上：\texttt{Impact} 最小 & LCIF 的在线近似，用历史表长度限界 \\
历史表 & 每 CPU 双向链表，上限 512 条 & 超限裁 1/4 并在锁外 \texttt{kfree}，保证在线性 \\
限流 & \texttt{YAT\_DECISION\_\allowbreak MIN\_INTERVAL\_NS} & 把决策频率写成常量（1\,ms），而不是靠负载碰巧 \\
\bottomrule
\end{tabular}
\end{table}

三处取舍值得说明：

\textbf{用时间近似重用距离。}严格的复用距离要逐次访问记账，代价是每次内存访问都要进内核路径。
内核改为累计"该核上与该作业相关的执行区间"，用时间差 $t(u_{j,v})-t(u_{j,v-1})$ 近似——
误差来自区间内其它作业的干扰，但省掉了逐访问钩子。

\textbf{候选集先按可运行性裁剪。}MSF 只在 \texttt{cpus\_allowed} 且当前可运行的核上计算，
这既是正确性要求（不能把任务放到不可运行的核），也是性能要求（避免 $O(m)$ 之外再乘历史表长度）。

\textbf{决策限流是常量而非启发式。}两次全量选核之间至少隔 1\,ms。这条不写在论文的算法伪代码里，
但它决定了实测中 \texttt{balance\_runs} 的量级（第\ref{sec:test}节），因此必须作为工程约束显式列出。

\subsection{复杂度与工程约束}

\begin{itemize}
  \item 预测一次命中级别是 $O(1)$（定长环形记录），选核是 $O(m)$，其中 $m$ 是可运行核数。
  \item 内核侧的历史表是每个 CPU 一条双向链表，查询是 $O(n_r)$；$n_r$ 上限 \texttt{YAT\_HISTORY\_MAX\_RECORDS}=512，超过后一次裁掉 1/4 并在锁外 \texttt{kfree}。
  \item 决策频率由 \texttt{YAT\_DECISION\_MIN\_INTERVAL\_NS}=1\,ms 限流：两次全量选核之间至少隔 1\,ms，把"每任务每秒最多 1000 次决策"这件事写进常量而不是靠负载碰巧。
  \item 一轮 balance 最多分配 \texttt{YAT\_BALANCE\_MAX\_BATCH}=256 个作业，这是防死循环的上限，不是吞吐目标。
\end{itemize}